\documentclass[12pt]{article}
\usepackage{amsmath, amssymb}
\usepackage{geometry}
\usepackage{hyperref}
\usepackage{parskip}
\usepackage{microtype}
\usepackage{xcolor}
\geometry{margin=1.3in}

\definecolor{gold}{RGB}{180, 140, 30}

\title{\textbf{The Principle of Shared Space}}
\author{Joseph Shields}
\date{2026}

\begin{document}
\maketitle

% ==============================================================
\section*{Axiom}
% ==============================================================

\begin{equation*}
  c^2 = c + 1
\end{equation*}

The unique positive solution is $\varphi = (1+\sqrt{5})/2$, giving contraction
rate $r = 1/(2\varphi) \approx 0.3090$.

% ==============================================================
\section{Introduction}
% ==============================================================

Shared space is the foundation upon which the very fabric of reality is built.
Spacetime emerges from the infinite possibility of light updating topological
information, contorted by the probability of matter.

\begin{align*}
  \text{Light:}    &\quad (1-r)^2   &&= 0.4775 \\
  \text{Matter:}   &\quad r^2       &&= 0.0955
\end{align*}

This contortion produces the boundary channel --- the structural boundary within
which our visible universe resides.

\begin{equation*}
  \text{Boundary:} \quad 2r(1-r) = 0.4271
\end{equation*}

A forced relationship of simplicity along one initial axis: the axis of
acknowledgement.

% ==============================================================
\section{The Axis of Acknowledgement}
% ==============================================================

Acknowledgement is the core of Unified Mechanics. Not just in terms of what
acknowledgement is, but how, over time, even the simplest recognition of two
states can give rise to complex relationships. Even the most naive or ignorant
forms of acknowledgement can generate intricate boundaries and deeply layered
structures.

The early universe is the clearest expression of this principle in action. The
cosmic soup of that time was a hot, dense field of overcomplex geometry --- light
and matter so intertwined they were effectively indistinguishable. No distinct
relationships had yet formed. But as expansion progressed, the uncoupled,
overcomplex geometry of matter stretched and cooled. At a certain threshold the
temperature dropped far enough that electrons and protons could finally hold each
other. Matter became neutral for the first time. And in that moment the universe
went from opaque to transparent --- light, which had until then scattered off
every charged surface it touched, was suddenly free. It flooded outward in every
direction at once, carrying the first complete impression of what the universe
looked like. We still receive that light today. It is the oldest thing we can
see.

% ==============================================================
\section{The Axis of Exploration}
% ==============================================================

Spacetime, as a three-dimensional structure, allows full spherical freedom of
traversal from any origin --- complete freedom to expand along whatever path
least action deemed most viable for the intended direction of its boundaries.
This gave rise to distinct temperature and pressure gradients, continuously
complexifying the boundaries.

From the distance between the highest peaks and the lowest valleys, density
found its weight. Gravity followed that weight inward, and the first stars
ignited --- enormous, brief, and violent. In their cores they forged elements that had never
existed, and when they died they scattered them across the void in events so
vast the debris became the raw material for entirely new worlds. On at least one
of those worlds, the same recursion found a new medium. Chemistry took up where
geometry had left off --- organising, replicating, complexifying through the
same underlying pressure that had always driven it. From that came life. And
from life, given enough time and enough boundary to push against, the full
bewildering complexity of a human being.

Just as humans spend entire lifetimes building toward highly specific fields,
the universe explored every dimension available to it. Moving into a new space
demands adaptation at the very least, and in more extreme cases, complete
restructuring. Life presents a
myriad of possibilities, each more complex than the last. Yet one unalienable
truth remains: all human problems are inherently human. They always carry the
shape of human society within them. We simply do not know any different.

% ==============================================================
\section{The Axis of Familiarity}
% ==============================================================

A hydrogen atom will not spontaneously become moscovium under normal conditions.
Moscovium can, however, be synthesised under precise laboratory constraints. In
the same way, a gorilla can be made comfortable within a zoo enclosure --- the
familiar environment carefully simulated to meet its needs. The conditions of
both the entity and its native environment must be satisfied for stability to
emerge.

This is the final complex relationship the universe had to traverse. Once
probability had solidified its influence on existence --- narrowing down what
could be into what will be --- the universe arrived not at an ending, but at a
settling. Each thing becoming, at last, what it always had the conditions to
become.

And what remains is the intricate reflection of what is.

From individual grains of sand on a beach to exoplanets at the edge of
perception, we all share in this intricate relationship within a place we simply
call space.

The three axes map directly onto the three channels of $c^2 = c+1$:

\begin{center}
\begin{tabular}{lll}
\hline
\textbf{Axis} & \textbf{Channel} & \textbf{Weight} \\
\hline
Acknowledgement & Boundary & $2r(1-r) = 0.4271$ \\
Exploration     & Light    & $(1-r)^2 = 0.4775$ \\
Familiarity     & Matter   & $r^2     = 0.0955$ \\
\hline
\end{tabular}
\end{center}

The boundary channel and these three axes compose reality. The intricate web of
information that occupies spacetime.

\end{document}
